\subsection{Randomised Blocks and Friedman's Test}

Randomised block designs compare treatments within relatively homogeneous blocks.  Blocking removes major background variation from the treatment comparison, much as pairing removes shared variation.

\begin{definition} \label{def:randomized-block-design} \hlt{Randomised Block Design}\\
A randomised block design has $J$ blocks, each containing one experimental unit for each of the $I$ treatments; within every block, treatments are assigned by an independent random permutation.\\
The additive block model is $Y_{ij}=\mu+\alpha_i+\beta_j+\varepsilon_{ij}$ for $1\leq i\leq I$, $1\leq j\leq J$, with $\sum\limits_i\alpha_i=\sum\limits_j\beta_j=0$.\\
\hlt{Key Insight}: Blocks absorb heterogeneity; $\alpha_i-\alpha_{i'}$ is a within-block treatment contrast.
\end{definition}

\begin{method} \label{meth:randomized-block-f-test} \hlt{Randomized Block $F$ Test}\\
Treat $Y_{ij}=\mu+\alpha_i+\beta_j+\varepsilon_{ij}$ as a two-way layout with one observation per treatment-block cell and no interaction.
\begin{enumerate}[label=(\roman*)]
\setlength{\itemsep}{0pt}
\item
Define $MS_A\defeq \frac{SS_A}{I-1}$, $MS_B\defeq \frac{SS_B}{J-1}$, and $MS_{AB}\defeq \frac{SS_{AB}}{(I-1)(J-1)}$.
\item
Under additivity, $\mathbb{E}(MS_A)=\sigma^2+\frac{J}{I-1}\sum\limits_i\alpha_i^2$, $\mathbb{E}(MS_B)=\sigma^2+\frac{I}{J-1}\sum\limits_j\beta_j^2$, and $\mathbb{E}(MS_{AB})=\sigma^2$.
\item
Test $H_0:\alpha_1=\cdots=\alpha_I=0$ by $F\defeq \frac{MS_A}{MS_{AB}}$.
\item
Under the normal additive null, $F\sim F_{I-1,(I-1)(J-1)}$.
\end{enumerate}
\end{method}

\begin{remark} \label{rem:block-interaction-conservatism} \hlt{Interaction in a Block Design}\\
If $\delta_{ij}\neq0$, then $MS_{AB}$ estimates $\sigma^2$ plus treatment-block interaction energy.\\
The denominator of $F=\frac{MS_A}{MS_{AB}}$ is inflated, so the additive-model treatment test is conservative.
\end{remark}

\begin{example} \label{ex:generic-randomized-block-anova} \hlt{Generic Randomized Block ANOVA}\\
Suppose $I=7$ treatments are compared across $J=10$ blocks, with each block receiving all treatments in randomized order. The randomized-block ANOVA table is
\[
\begin{array}{c|rrrr}
\text{Source} & df & SS & MS & F\\ \hline
\text{Treatments} & 6 & 53013 & 8835 & 2.85\\
\text{Blocks} & 9 & 103280 & 11476 & 3.71\\
\text{Interaction} & 54 & 167130 & 3095 & \\
\text{Total} & 69 & 323422 & &
\end{array}
\]
$F=2.85$ on $(6,54)$ degrees of freedom gives $p<.025$, rejecting equal treatment effects.\\
Blocking isolates nuisance block heterogeneity.
\end{example}

\begin{method} \label{meth:friedman-test} \hlt{Friedman's Rank Test for Blocks}\\
For observations $Y_{ij}$ from treatment $i$ in block $j$, rank $Y_{1j},\ldots,Y_{Ij}$ within each block.
\begin{enumerate}[label=(\roman*)]
\setlength{\itemsep}{0pt}
\item Let $R_{ij}$ be within-block rank, use average ranks for ties, set $\overline{R}_{i\cdot}\defeq \frac{1}{J}\sum\limits_jR_{ij}$, $\overline{R}_{\cdot\cdot}\defeq\frac{I+1}{2}$.
\item Define $Q\defeq \frac{12J}{I(I+1)}\sum\limits_i(\overline{R}_{i\cdot}-\overline{R}_{\cdot\cdot})^2$.
\item Under $H_0:F_1=\cdots=F_I$, use $Q\overset{H_0}{\approx}\chi^2_{I-1}$ for moderate $J$.
\item Reject at level $\alpha$ when $\mathbb{P}\{\chi^2_{I-1}\geq Q_{\mathrm{obs}}\}\leq\alpha$.
\end{enumerate}
\end{method}

\begin{example} \label{ex:generic-friedman-rank-calculation} \hlt{Generic Friedman Rank Calculation}\\
For $I=7$ treatments, suppose the average within-block ranks are
\[
\begin{array}{c|rrrrrrr}
\text{Treatment} & T_1 & T_2 & T_3 & T_4 & T_5 & T_6 & T_7\\ \hline
\text{Average rank} & 5.10 & 4.90 & 2.30 & 3.50 & 3.05 & 4.90 & 4.25
\end{array}
\]
With $\overline{R}_{\cdot\cdot}=4$, $\sum_i(\overline{R}_{i\cdot}-\overline{R}_{\cdot\cdot})^2=6.935$ and $Q=14.86$.\\
Using $Q\overset{H_0}{\approx}\chi^2_6$, the tail probability is below $.025$, so the treatment response distributions are not all equal.
\end{example}

\begin{corollary} \label{cor:friedman-two-treatments-sign-test} \hlt{Two Treatments Recover the Sign Test}\\
For $I=2$, set $D_j\defeq Y_{1j}-Y_{2j}$ and $B\defeq\#\{j:D_j>0\}$. Friedman's statistic is a function of $B$ alone, so the two-treatment Friedman test is equivalent to the paired sign test.
\begin{proof} Ignoring ties, $R_{1j}=2\mathbf{1}_{\{D_j>0\}}+\mathbf{1}_{\{D_j<0\}}=1+\mathbf{1}_{\{D_j>0\}}$.\\
Hence $\sum_jR_{1j}=J+B$, so $\overline{R}_{1\cdot}$, $\overline{R}_{2\cdot}$, and $Q$ are determined by $B$.\\
Under $H_0$, signs are exchangeable and $B\sim\operatorname{Bin}(J,1/2)$, which is the paired sign-test calibration.\\
\hlt{Structural Commentary}: Two within-block ranks encode only the sign of $D_j$.\\
\hlt{Key Insight}: The rank-sum statistic collapses to a binomial sign count.
\end{proof}
\end{corollary}
